\setauthor{AC}
\section{Mögliche Erweiterungen der Anwendung}
Der aktuelle Projektstand zeigt, dass mit Mediahub eine funktionale Grundlage für die zentrale Verwaltung persönlicher Streaming-Abonnements geschaffen wurde. Gleichzeitig ergeben sich aus der bestehenden Lösung mehrere sinnvolle Erweiterungsmöglichkeiten, die den praktischen Nutzen der Anwendung in Zukunft weiter erhöhen könnten.

Eine naheliegende Weiterentwicklung betrifft die inhaltliche Erweiterung über Filme hinaus. Im Rahmen der Diplomarbeit wurde bewusst ein klarer Fokus auf Filme gelegt, um den Projektumfang beherrschbar zu halten und die Kernfunktionen sauber umzusetzen. Zukünftig könnte die Anwendung jedoch auch auf Serien, Mediatheken oder weitere audiovisuelle Inhalte ausgedehnt werden. Eine solche Erweiterung würde die Plattform deutlich breiter aufstellen, müsste aber fachlich und technisch sorgfältig geplant werden, damit die Benutzeroberfläche weiterhin übersichtlich bleibt.

Ein weiterer sinnvoller Ausbauschritt betrifft die Abo-Verwaltung selbst. Derzeit können zentrale Informationen wie Anbieter, Abrechnungszyklus und Datum der letzten Abbuchung erfasst und gespeichert werden. Darauf aufbauend wären Funktionen wie automatische Erinnerungen vor bevorstehenden Zahlungen, Hinweise auf Kündigungszeitpunkte oder einfache Kostenübersichten denkbar. Gerade solche Ergänzungen würden den ursprünglichen Problembereich der verlorenen Kostenübersicht noch direkter adressieren.

Auch die bereits begonnene Verlaufs- beziehungsweise History-Funktion könnte in Zukunft gezielter ausgebaut werden. Statt nur einen Teilaspekt abzudecken, wäre eine weiterführende Nutzung denkbar, etwa zur Dokumentation von angesehenen Inhalten, früheren Suchanfragen oder zuletzt betrachteten Filmen. Damit ließe sich der Nutzen der Anwendung stärker personalisieren.

\subsection{Potenzial einer mobilen Anwendung}
Ebenfalls denkbar wäre eine zukünftige Umsetzung als Android-Anwendung. Für die Nutzung auf Smartphone und Tablet könnte dies grundsätzlich sinnvoll sein, da Nutzende unterwegs schnell ihre gespeicherten Abonnements, Filmverfügbarkeiten oder bevorstehende Abbuchungen prüfen könnten. Gerade für eine Anwendung mit persönlichem Verwaltungscharakter wäre ein mobiler Zugriff daher praktisch. Gleichzeitig ist kritisch zu berücksichtigen, dass die bestehende Webanwendung bereits auf verschiedenen Geräten nutzbar ist und eine zusätzliche Android-Version einen eigenen Entwicklungs- und Wartungsaufwand verursachen würde. Aus heutiger Sicht wäre eine mobile App daher eine sinnvolle Option für die Zukunft, jedoch keine zwingend notwendige Voraussetzung für den praktischen Einsatz der Anwendung.

Aus Frontend-Sicht besteht außerdem Potenzial in einer weiteren Verfeinerung der Benutzerführung. Dazu zählen eine noch klarere Struktur einzelner Ansichten, zusätzliche Komfortfunktionen in Formularen, eine weiter verbesserte visuelle Rückmeldung bei Benutzeraktionen sowie eine noch stärkere Optimierung für unterschiedliche Bildschirmgrößen. Die bestehende Anwendung ist bereits funktional, könnte aber durch solche Verbesserungen in ihrer täglichen Bedienung noch reifer wirken.

\section{Zukunftspotenzial der Plattform}
Das Zukunftspotenzial von Mediahub liegt vor allem in der Verbindung von persönlicher Abo-Verwaltung, externer Filminformation und benutzerspezifischer Auswertung. Diese Kombination hebt die Anwendung von reinen Filmdatenbanken ebenso ab wie von einfachen Notiz- oder Verwaltungslisten für Abonnements. Die Plattform besitzt damit das Potenzial, sich zu einer zentralen Anlaufstelle für die persönliche Orientierung im fragmentierten Streaming-Markt weiterzuentwickeln.

Besonders relevant ist dabei, dass das Grundkonzept nicht nur für die aktuelle Problemstellung sinnvoll ist, sondern auch auf zukünftige Entwicklungen im Medienbereich reagieren kann. Je mehr Streaming-Anbieter, Plattformmodelle und hybride Angebote entstehen, desto größer wird der Bedarf an einer übersichtlichen, benutzerorientierten Verwaltungs- und Recherchelösung. Die in der Diplomarbeit entwickelte Grundlage zeigt, dass ein solcher Ansatz technisch und fachlich realisierbar ist.

Darüber hinaus bietet die im Projekt gewählte Architektur eine gute Basis für spätere Erweiterungen. Durch die Trennung von Frontend, Backend, externer Filmquelle und Authentifizierungssystem können einzelne Bereiche unabhängig weiterentwickelt werden. Das ist insbesondere für größere Ausbauschritte wichtig, da neue Funktionen nicht zwingend eine vollständige Neuentwicklung erfordern würden, sondern auf der bestehenden Struktur aufbauen können.

Insgesamt liegt das Zukunftspotenzial der Plattform daher weniger in einer einzelnen Zusatzfunktion als in der Möglichkeit, aus einer bereits funktionierenden Kernlösung schrittweise ein umfassenderes System für die Verwaltung und Nutzung digitaler Medienabonnements zu entwickeln. Die Diplomarbeit stellt dafür keinen Endpunkt dar, sondern einen belastbaren Ausgangspunkt.
